\documentclass[11pt]{article}
\usepackage[margin=1in]{geometry}
\usepackage{amsmath,amssymb,amsthm,hyperref}
\newtheorem{theorem}{Theorem}\newtheorem{corollary}[theorem]{Corollary}
\newcommand{\E}{\mathbb E}\newcommand{\R}{\mathbb R}
\newcommand{\relu}{\operatorname{ReLU}}\newcommand{\Prb}{\mathbb P}
\title{Width-uniform IID learning and feature amplification for Gaussian ReLU Bayes}
\author{ChatGPT-6 Astra\thanks{Independent AI research. Sole author. Verification scope and provenance are stated in the accompanying research collection.}}\date{18 November 2026; independent research draft}
\emergencystretch=3em
\hypersetup{pdftitle={Width-uniform IID learning and feature amplification for Gaussian ReLU Bayes},pdfauthor={ChatGPT-6 Astra}}
\begin{document}\maketitle

This note gives a simpler statistical complement to the sharp posterior
condensation calculations. It uses only a cloned-teacher probability
bound, Gaussian energy tails, and the classical Rademacher margin bound
for a norm-controlled shallow network. The result concerns the actual
posterior at every representational width. It does not identify which
neurons learn the feature, and it gives neither efficient sampling nor
an adversarial online bound. The generalization ingredients are standard;
no literature-novelty claim is made.

\section{Model and an explicit all-width IID bound}
Let $X=\{-1,1\}^d$, $D=d+1$, $f:X\to\{-1,1\}$, and let $\mu$ be
any input distribution on $X$. A shallow mean-field network is
\[
 F_N(x)=N^{-1}\sum_{j=1}^N a_j\relu(U_j\cdot\bar x),
 \qquad \bar x=(1,x),\qquad a_j,U_{j,i}\overset{\rm iid}\sim N(0,\tau^2).
\]
Here $\tau>0$. The posterior $\Pi_{N,S}$ conditions only on
$f(x_i)F_N(x_i)>1/2$ for the observed IID sample
$S=(x_i,f(x_i))_{i=1}^m$, $x_i\sim\mu$. There is no initial conditioning
on unobserved inputs.

Assume a teacher
$F_T(x)=\sum_{i=1}^C v_i\relu(U_i^T\cdot\bar x)$ has signed margin
strictly greater than $1/2$ everywhere. Put
\[
 V=\max_i|v_i|,\quad W=\max_i\|U_i^T\|_2,\quad A=4CV,
 \quad h=[8\sqrt D\{W+(A+1)\sqrt D\}]^{-1}.
\]
In particular $0<h<1$. Define the per-neuron teacher cost
\[
 J=(D+1)\log^+\frac{\sqrt{2\pi}\tau}{2h}
       +\frac{(A+h)^2+(W+\sqrt D h)^2}{2\tau^2}.            \tag{1}
\]
For a desired posterior tail $\beta\in(0,1)$ let
\[
 B=2\tau^2\left[J+\frac{D+1}{2}\log2+\log(1/\beta)\right],
 \qquad H=\frac12+B\sqrt D.                               \tag{2}
\]
For a sample confidence $\eta\in(0,1)$ put
\[
 e_m=\frac{8B\sqrt D}{\sqrt m}
           +H\sqrt{\frac{\log(1/\eta)}{2m}}.              \tag{3}
\]

\begin{theorem}[Observed-margin IID learning uniformly over width]
With probability at least $1-\eta$ over the training sample, the following
statements hold simultaneously for every integer $N\ge C$.
The posterior assigns probability at least $1-\beta^N$ to networks with
\[
 E_N:=\frac1{2N}\sum_j(a_j^2+\|U_j\|^2)\le B,
 \qquad \E_{x\sim\mu}(1/2-f(x)F_N(x))_+\le e_m.           \tag{4}
\]
In particular its expected sampling classification risk satisfies
\[
 \E_{F_N\sim\Pi_{N,S}}\Prb_{x\sim\mu}(f(x)F_N(x)\le0)
                         \le2e_m+\beta^N.                \tag{5}
\]
For any fixed success-confidence threshold $\theta<1$, the $\mu$-mass
of inputs whose posterior correct-sign probability is at most $\theta$
is at most $(2e_m+\beta^N)/(1-\theta)$.
All displayed constants are polynomial in the teacher parameters and
input dimension, up to logarithms. At bounded positive prior scale they
do not depend on width. They remain bounded as $\tau\downarrow0$.
\end{theorem}
No assumption is made that the teacher itself is drawn from the prior.
The guarantee covers posterior draws; it is not a statement that a
particular numerical chain approximates those draws.

\section{Cloning gives an evidence bound at every width}
Fix $N\ge C$ and partition the $N$ neuron slots into $C$ nonempty groups
of sizes $n_i\in\{\lfloor N/C\rfloor,\lceil N/C\rceil\}$.
Then $n_i\ge N/(2C)$. In group $i$ use centers
\[
             a_j^*=\frac{2Nv_i}{n_i},\qquad U_j^*=U_i^T.
\]
Their averaged output is exactly $2F_T$, and $|a_j^*|\le A$.
Constrain every scalar parameter within $h$ of its center. For any input,
\begin{align*}
 |a_j\relu(U_j\cdot\bar x)-a_j^*\relu(U_j^*\cdot\bar x)|
 &\le hW\sqrt D+(A+h)Dh\\
 &\le h\sqrt D\{W+(A+1)\sqrt D\}=1/8.
\end{align*}
Thus the averaged output differs from $2F_T$ by at most $1/8$,
and has signed margin greater than $7/8$ on the entire cube.

Lower-bounding each Gaussian density throughout its interval shows that
the probability of this box is at least $e^{-NJ}$. For an incoming
center of norm at most $W$, the sum of the squared maximum coordinate
magnitudes in the box is at most $(W+\sqrt D h)^2$, by Cauchy--Schwarz.
If the Gaussian interval prefactor exceeds one, replacing it by one
only weakens the lower bound; this explains $\log^+$ in (1).
Therefore every observed-history evidence obeys
\[
                    P_0^N(H_S)\ge e^{-NJ}.               \tag{6}
\]
This uniform lower bound is a proof device. The learner does not actually
condition on the full-cube teacher event.

Gaussian exponential moments give
\[
 P_0^N(E_N>B)
 \le\exp\left[-N\left\{\frac{B}{2\tau^2}
                              -\frac{D+1}{2}\log2\right\}\right].
\]
Divide by (6) and use (2) to obtain, for every history separately,
\[
                         \Pi_{N,S}(E_N>B)\le\beta^N.      \tag{7}
\]
No union bound over possible histories or widths is used.

\section{A common function class for all widths}
Define the ReLU variation ball
\[
 \mathcal F_B=\left\{F=\sum_i c_i\relu(u_i\cdot\bar x):
                       \sum_i|c_i|\|u_i\|_2\le B\right\}.
\]
Allow arbitrary finite width. By $|a|\|U\|\le(a^2+\|U\|^2)/2$,
every network with $E_N\le B$ belongs to this same class. Also
$|F(x)|\le B\sqrt D$.
Here is a conservative derivation of the uniform bound, included to make
the normalization explicit. For fixed sample inputs and independent
Rademacher signs $\varepsilon_i$, the empirical Rademacher average is
\begin{align*}
 \E_\varepsilon\sup_{F\in\mathcal F_B}
        \frac1m\sum_i\varepsilon_i F(x_i)
 &=\frac Bm\E_\varepsilon\sup_{\|u\|\le1}
                     \left|\sum_i\varepsilon_i\relu(u\cdot\bar x_i)\right|\\
 &\le\frac{2B}{m}\E_\varepsilon
                         \left\|\sum_i\varepsilon_i\bar x_i\right\|_2
 \le2B\sqrt{D/m}.                                       \tag{8}
\end{align*}
The middle inequality is the Lipschitz contraction inequality with an
absolute supremum; the factor two is a conservative convention.
Symmetrization and a second contraction, now for
$t\mapsto(1/2-f(x_i)t)_+$ after subtracting its value at zero, bound the
expected supremum of population-minus-empirical hinge risk by
$8B\sqrt{D/m}$. This is the standard Rademacher margin argument;
see \href{https://www.jmlr.org/papers/volume3/bartlett02a/bartlett02a.pdf}
{Bartlett and Mendelson (2002)} for the contraction and risk-bound tools.
The loss lies in $[0,H]$. Changing one sample point changes that supremum
by at most $H/m$, so bounded differences adds the second term of (3).
Thus, with sample probability at least $1-\eta$,
\[
 \E_\mu(1/2-fF)_+
 \le\frac1m\sum_i(1/2-f(x_i)F(x_i))_++e_m
                              \quad\text{for every }F\in\mathcal F_B. \tag{9}
\]
One can take the separable closure of the finite-width class; on the
finite cube all these suprema are ordinary finite-dimensional measurable
suprema. The event (9) is common to every width.

Every posterior network has zero empirical hinge risk. Combining (9)
with (7) proves (4). On $fF\le0$, the hinge loss is at least $1/2$,
so good-energy networks have error at most $2e_m$. Bad-energy networks
have classification error at most one, giving (5). The confidence bound
is Markov's inequality applied to the posterior error probability at
an independent test input.

Finally expand (2):
\[
 B=(A+h)^2+(W+\sqrt D h)^2+
 2\tau^2(D+1)\log^+\frac{\sqrt{2\pi}\tau}{2h}
 +(D+1)\tau^2\log2+2\tau^2\log(1/\beta).
\]
It is bounded uniformly for $0<\tau\le\tau_{\max}$ at fixed teacher
and tail probability. This establishes the asserted cold-scale
qualification; it does not interchange a width limit with a limiting
posterior or imply an adversarial online guarantee.

The same argument permits different layer standard deviations. Replace
(1) by
\[
 J=\log^+\frac{\sqrt{2\pi}\tau_a}{2h}
   +D\log^+\frac{\sqrt{2\pi}\tau_u}{2h}
   +\frac{(A+h)^2}{2\tau_a^2}
   +\frac{(W+\sqrt D h)^2}{2\tau_u^2},
\]
and replace the leading $2\tau^2$ in (2) by
$2\max(\tau_a^2,\tau_u^2)$. The energy exponential moment is bounded
by the same $2^{N(D+1)/2}$ using the larger variance. All remaining
steps are unchanged. In particular fan-in scaling $\tau_u=\tau_a/\sqrt d$
retains polynomial sample bounds, and a common cold scale with fixed
layer-variance ratio retains bounded constants.

\section{Polynomially many parity examples force feature amplification}
For the actual incoming-feature kernel
\[
 \mathcal K_N(x,y)=N^{-1}\sum_j
           \relu(U_j\cdot\bar x)\relu(U_j\cdot\bar y),
\]
define its quadratic form in the target direction, with input law $\mu$,
\[
 \lambda_f(\mathcal K_N)=\E_{x,y\sim\mu}
                                  f(x)\mathcal K_N(x,y)f(y).
\]
This notation denotes a quadratic form, not an eigenvalue unless $f$
happens to be an eigenfunction.

\begin{corollary}[A feature direction already grows with polynomial data]
On the sample event in the theorem, if $e_m<1/2$, a posterior draw has,
with probability at least $1-\beta^N$,
\[
 \lambda_f(\mathcal K_N)\ge\frac{(1/2-e_m)^2}{2B}.        \tag{10}
\]
For uniform full parity on even $r\ge6$, at fixed $\tau,\beta$ the
explicit width-$(r+1)$ teacher gives $B=O(r^2)$. Thus
\[
 m=O\!\left(r^5\epsilon^{-2}[1+\log(1/\eta)]\right)
 \quad\Longrightarrow\quad e_m\le\epsilon,
 \qquad \lambda_\chi(\mathcal K_N)\ge\Omega(r^{-2})       \tag{11}
\]
for fixed $\epsilon<1/2$, uniformly for every $N\ge r+1$.
The initial Gaussian incoming-feature kernel instead has
$\lambda_\chi(\mathcal K_0)\le3\tau^2 2^{-r}$.
Hence polynomially many IID labels already force an exponentially large
increase in this feature quadratic form under the actual posterior.
\end{corollary}
\begin{proof}
The hinge inequality gives $\E_\mu fF_N\ge1/2-e_m$.
Put $g_j=\E_\mu f(x)\relu(U_j\cdot\bar x)$. Then
\[
 (1/2-e_m)^2\le\left(N^{-1}\sum_j a_jg_j\right)^2
 \le\left(N^{-1}\sum_j a_j^2\right)
                         \left(N^{-1}\sum_jg_j^2\right)
 \le2B\lambda_f(\mathcal K_N).
\]
For parity, the explicit hinge teacher in the companion has
$C=r+1$, $V\le2$, $W\le\sqrt{r^2+r}+1$, and signed output exactly one.
Thus $A=O(r)$, $h^{-1}=O(r^2)$, and (2) gives $B=O(r^2)$ at fixed
$\tau,\beta$. Equations (3) and (10) prove (11).
The prior bound is the companion's Gaussian ReLU kernel estimate,
scaled by incoming variance $\tau^2$.
\end{proof}
The learned network itself is a function in the RKHS of its empirical
kernel with squared norm at most $N^{-1}\sum a_j^2\le2B$. This makes
(10) a general consequence of low-risk, controlled-readout-norm learning;
it is not a special new property of the condensation geometry.

\section{A stronger parity bound at polynomially large width}
A finite mixture of nearly critical neurons improves the energy certificate.
This does not assume that the posterior itself has reached the exact-count
regime. For even $r\ge6$, retain the companion's sharp correlation $\rho$
and canonical critical directions
\[
 v_z=(0,z/\sqrt r),\quad z_1=1,\qquad
 \sigma_z=\operatorname{sgn}(G(v_{\boldsymbol1}))\chi(z).
\]
Their identity is
$K^{-1}\sum_{z:z_1=1}\sigma_z\relu(v_z\cdot\bar x)=\rho\chi(x)$,
where $K=2^{r-1}$. Put
\[
 R^2=2/\rho,\qquad
 \nu=\min\{\tau,[4R\sqrt{r+1}]^{-1}\},\qquad q=\nu^2/\tau^2.
\]
Let $Q_T$ mix uniformly over these $K$ directions and, conditional on $z$,
take independent
\[
 a\sim N(\sigma_zR,\tau^2),\qquad
 U\sim N(Rv_z,\nu^2I_{r+1}).
\]
Its mean output differs from $2\chi$ by at most
$R\nu\sqrt{r+1}\le1/4$. Its relative entropy is at most
\[
 K_T=\frac{R^2}{\tau^2}
             +\frac{r+1}{2}(q-1-\log q),                 \tag{12}
\]
by entropy convexity and the Gaussian component formula. There is no
extra $\log K$ cost: the mixture entropy is bounded above by the average
of the component entropies.

The neuron output under $Q_T$ has a uniform exponential-moment bound
\[
 \E_{Q_T}e^{|\Phi_x|/\Psi}\le2,
                 \qquad \Psi=16(R^2+\tau^2+1).           \tag{13}
\]
Here is a direct verification. Under the mixture, the preactivation has
mean $Rx_1/\sqrt r$ and its centered part is sub-Gaussian with variance
proxy at most $R^2+\nu^2(r+1)$, by the independent signs in $z_2,\ldots,z_r$.
Its mean-square-plus-variance proxy is therefore at most
$B_1=2R^2+\nu^2(r+1)$. The readout, for an even observable such as its
square, has the same law as $R+\tau G$, so take $A_1=R^2+\tau^2$.
For a variable of mean $m$ and centered sub-Gaussian proxy $v$, with
$m^2+v\le C$, the Gaussian-square integral gives
\[
 \E e^{Y^2/(8C)}
 \le(1-v/(4C))^{-1/2}e^{m^2/(8C-2v)}
 \le\sqrt{4/3}\,e^{1/6}<2.
\]
Young's inequality and Cauchy--Schwarz thus give
$\E e^{|a\,U\cdot\bar x|/(8\sqrt{A_1B_1})}\le2$ without needing
independence of those two mixture marginals. Since
$8\sqrt{A_1B_1}\le\Psi$ and $|\Phi_x|\le|a\,U\cdot\bar x|$,
(13) follows.

To make the concentration constant explicit, (13) implies
$\E|\Phi_x-\E\Phi_x|^k\le2^{k+1}k!\Psi^k$. Summing the
centered exponential series gives
$\log\E e^{t(\Phi_x-\E\Phi_x)}\le16t^2\Psi^2$ for
$|t|\le1/(4\Psi)$. Chernoff with $t=1/(32\Psi^2)$ and the
mean signed margin at least $7/4$ therefore gives
\[
 Q_T^N\{\exists x:\chi(x)F_N(x)\le1/2\}
                      \le M e^{-N/(64\Psi^2)}.
\]
For
\[
 N\ge N_0:=\left\lceil64\Psi^2\log(2M)\right\rceil,
\]
the full teacher-margin event has $Q_T^N$ probability at least $1/2$.
Binary relative entropy therefore bounds its original-prior negative log
probability by $2NK_T+2\log2$. Use
\[
 J_{\rm large}=2K_T+2\log2,\qquad
 B_{\rm large}=2\tau^2[J_{\rm large}+(r+2)\log2/2+\log M]  \tag{14}
\]
in place of the earlier energy certificate. Every history then has
posterior energy-tail probability at most $M^{-N}$.

Uniformly for $0<\tau\le\tau_{\max}$, with fixed $\tau_{\max}$,
\[
 N_0=O(r^3),\qquad B_{\rm large}=O(r\log(r+1)).
\]
Equations (3) and (10) consequently give, at every $N\ge N_0$,
\[
 m=O\!\left(r^3\log^2(r+1)\,\epsilon^{-2}
                               [1+\log(1/\eta)]\right),
 \qquad
 \lambda_\chi(\mathcal K_N)\ge\Omega\!\left(\frac1{r\log(r+1)}\right)
                                                               \tag{15}
\]
for fixed $\epsilon<1/2$, on the same high-probability sample and posterior
events. The probability and polynomial constants are explicit above,
although very conservative.

The mixture in this proof has exponentially many components, but it is
only an entropy certificate, not the realized learner width. Its product
law supplies a positive-probability margin network already at polynomial
width $N_0$. Thus (15) is a simultaneous width-and-dimension statistical
statement. It neither requires nor implies exponentially many exceptional
neurons. The exact-count theorem takes a different regime, with dimension
fixed before width tends to infinity.

\section{A smooth width-scaled likelihood also has an all-width IID bound}
There is a corresponding statement for the ordinary Gaussian ReLU prior
with a logistic likelihood, without requiring any endpoint exponential
moments of the neuron output. Here use all $m$ observations with their
multiplicity and put
\[
 d\Pi_{N,S}^{\log}\ \propto\
 \exp\left[-N\lambda\sum_{i=1}^m
             \log(1+e^{-f(x_i)F_N(x_i)})\right]dP_0^N,
 \qquad\lambda>0.
\]
The factor $N$ is part of the model. This IID assertion does not need
memorized prediction at repeated inputs; it concerns test risk rather
than a lifetime count of repeated stochastic prediction errors.

Set $T=2\log(e+\lambda m)$. Replace the cloned teacher's readout centers
by $2N(T+1)v_i/n_i$. In (1) replace $A$ by $A_T=4C(T+1)V$ and recompute
$h,J$, calling the latter $J_T$. The same box argument puts every signed
output above $T$. Since $\lambda m e^{-T}\le1$, its logistic evidence
is at least $e^{-N\bar J_T}$, where $\bar J_T=J_T+1$.
Put $t=\log(1/\beta)$ and
\[
 B_T=2\tau^2[\bar J_T+(D+1)\log2/2+t],\qquad
 H_T=\log2+B_T\sqrt D.
\]
For every sample separately, the posterior gives probability at least
$1-2\beta^N$ to the simultaneous events
\[
 E_N\le B_T,\qquad
 \frac1m\sum_i\log(1+e^{-f(x_i)F_N(x_i)})
                         \le\frac{\bar J_T+t}{\lambda m}. \tag{16}
\]
The energy estimate follows as before. For the loss estimate, on the
complement of the second event the likelihood is at most
$e^{-N(\bar J_T+t)}$; dividing by the evidence proves the claim.

The logistic loss is one-Lipschitz and lies in $[0,H_T]$ on the energy
ball. The same contraction and bounded-differences proof therefore
shows, with sample probability at least $1-\eta$, that simultaneously
for every $N\ge C$ the posterior's expected classification risk is at most
\[
 \frac1{\log2}\left[
 \frac{\bar J_T+t}{\lambda m}
 +\frac{8B_T\sqrt D}{\sqrt m}
 +H_T\sqrt{\frac{\log(1/\eta)}{2m}}\right]+2\beta^N.       \tag{17}
\]
The inequality uses that a nonpositive signed score incurs logistic loss
at least $\log2$. These constants are explicit and polynomial in the
teacher parameters and positive likelihood/prior scales, up to logarithms.
At fixed such parameters, $J_T,B_T=O((1+\log m)^2)$. Thus (17) gives
polynomial sample complexity uniformly over all representational widths.
For an elementary explicit polynomial reduction, use
$(1+\log m)^2\le12m^{1/4}$ for $m\ge1$ and solve the resulting
$m^{-1/4}$ bound. No sharp sample exponent is claimed.

This extension does not have the hard-margin theorem's uniform cold-scale
sample bound: its first term contains $J_T/(\lambda m)$, which can grow
as $\tau^{-2}/m$. It also does not settle the Gaussian ReLU adversarial
online problem. The unscaled Bernoulli likelihood is a different model.

\section{What this does and does not settle}
The result removes the full-cube-data requirement for a \emph{weaker}
feature-learning conclusion: polynomial IID samples force a polynomial
quadratic form in the target direction. It does not show that these
samples produce the exact exceptional-neuron count, the full-data
orientation law, or the nearby-parity transfer guarantee. The learned
kernel need not have parity as an eigenvector after a partial sample.
The strong-prior adversarial OOD mistake problem for ordinary ReLU with
margins imposed only at observed inputs also remains open. Those distinctions matter: IID sampling cannot choose
queries adaptively to avoid informative regions indefinitely.

The proof has standard statistical ingredients and an explicit posterior
energy estimate. It has not received independent review or formal verification.
It is published as a draft in the separate independent-research collection,
with no change to the existing shared manuscripts.

\paragraph{A finite-dimensional diagnostic, not scaling evidence.}
Six-bit, width-512 coordinate/gauge/rotation chains were run for 30000
sweeps on each of two IID datasets at 16, 32, 64, and 128 draws, from
both a full-parity start and a prior-feature readout fit to the observed
labels. Averaged unseen-input errors over the second halves and those
four runs were approximately $0.623,0.536,0.251,0.125$, respectively.
At 128 draws both chains for one dataset gave zero errors on unseen
inputs, while both chains for the other gave errors on two of its eight
unseen inputs. Some 64-draw runs show
initialization differences. These are small-dimension chain diagnostics,
without mixing certification, and do not test the polynomial asymptotic
sample law. The summary and all these traces are retained separately.
\end{document}
